拍拍贷成熟窗口保留73,696个成熟可判定首标，其中16,446例事件。训练/验证/测试分别为13,722/17,681/10,347笔，对应事件3,452/4,068/2,099例。训练放款截至2015年末，窗口最晚结束2016年3月；验证放款2016年4至6月，窗口最晚结束9月；测试放款10月1日至11月22日，窗口截至2017年2月22日。源表在统一回溯时点发布，时间切分采用逻辑成熟窗口；首标及零历史成功借款作为客户隔离代理。

标签依据到期日、实际还款日、状态及正金额还款义务共同构造；未到期、零义务、计划不完整及日期冲突保留未知。正金额义务条件排除了训练期4笔零应还却标迟延记录，以及209笔窗口内无正金额应还记录；验证和测试标签未变。按统一标签规则与固定配置重新计算全部模型。

\begin{table}[H]\centering\caption{拍拍贷：全部固定模型的测试表现}
\fontsize{9.1bp}{12.5bp}\selectfont
\begin{tabularx}{\linewidth}{@{}Yrrrrrr@{}}
\toprule
模型 & AUC & AP & KS & Logloss & Brier & 10\%覆盖 \\
\midrule
Logistic & 0.5483 & 0.2266 & 0.0760 & 0.5036 & 0.1615 & 259 \\
CatBoost & 0.5566 & 0.2423 & 0.0766 & 0.5006 & 0.1605 & 282 \\
LightGBM & 0.5429 & 0.2325 & 0.0561 & 0.5056 & 0.1621 & 267 \\
XGBoost & 0.5535 & 0.2371 & 0.0721 & 0.5013 & 0.1609 & 274 \\
三族等权集成 & 0.5524 & 0.2388 & 0.0690 & 0.5016 & 0.1609 & 274 \\
\bottomrule\end{tabularx}
\tabnote{两个阶段各拟合4个模型：开发选择及训练+验证最终拟合。等权集成不另算训练。逐行标签条件及模型配置见复算材料。}
\end{table}